\subsection{Technology Stack Selection}
\label{sec:tech_stack}

The technology stack for aBridgeAI was selected by evaluating each major system component against its most viable alternatives. Each decision is grounded in the system's core technical constraints: native AI/ML orchestration, real-time WebRTC voice processing, high-throughput vector search, and graph-based knowledge representation.

Each subsection below states the constraint the component must satisfy and the decision it led to. The comparison matrix behind each --- the chosen option scored against the alternatives on the dimensions that mattered --- is given in Appendix~\ref{appendix:techstack}.

% ─────────────────────────────────────────────
\subsubsection{Backend Framework}
% ─────────────────────────────────────────────

The backend must orchestrate adaptive SR scheduling, knowledge retention estimation, asynchronous LLM calls, and interview lifecycle operations. The implementation uses first-party Python modules for retrieval, state policy, and model integration, together with standard scientific libraries where appropriate; it does not depend on a chain-composition framework.


\textbf{Decision: FastAPI.} FastAPI combines Python-native AI library access with a fully asynchronous request model. It provides the lifecycle, authorization, token, dispatch, and persistence APIs for the native LiveKit interview worker, while independent worker processes hold long-running real-time sessions and asynchronous \texttt{arq} jobs. Its built-in OpenAPI generation also accelerates API contract development across the team.

% ─────────────────────────────────────────────
\subsubsection{Frontend Framework}
% ─────────────────────────────────────────────

The frontend must support both a standard LMS interface (course catalog, quiz editor, progress dashboards) and a real-time voice interview module with live transcript rendering and session state management. As a fully authenticated SPA with no public-facing SEO requirements, the system does not benefit from Server-Side Rendering (SSR) — making a lightweight CSR-first toolchain the optimal choice.


\textbf{Decision: Vite + React + TanStack Query.} aBridgeAI is a fully authenticated application — all routes require login, so SSR and SEO optimizations provided by Next.js yield no benefit. Next.js App Router introduces significant architectural complexity: React Server Components (RSC) require developers to reason about server/client boundaries throughout the codebase, the Flight payload serialization adds overhead at every server-client crossing, and the framework has accumulated a reputation for instability across major versions \cite{northflank2024nextjs, appwrite2024tanstack}. Vite produces a 54\% smaller initial bundle ($\sim$42KB vs $\sim$92KB) and delivers sub-50ms Hot Module Replacement, directly improving developer iteration speed \cite{techsy2026vite, hygraph2024vite}. TanStack Query replaces Next.js's coarse fetch caching with fine-grained server-state management — cache invalidation, background refetching, and optimistic updates — using standard React patterns that require no framework-specific mental model \cite{tanstack2024comparison}.

% ─────────────────────────────────────────────
\subsubsection{Real-Time Communication}
% ─────────────────────────────────────────────

The AI interview feature requires low-latency voice transport, streaming STT/TTS integration, and conversational agent turn-taking. The communication layer must handle continuous audio streams without TCP head-of-line blocking.


\textbf{Decision: LiveKit.} The deployed architecture uses LiveKit Cloud for the WebRTC room and a separate native LiveKit Agents worker for the interview runtime. The worker combines configured STT, LLM, and TTS plugins with Silero VAD and persistent conversational context; application code retains question selection and assessment state authority. Socket.io's TCP transport causes head-of-line blocking that makes continuous voice streams unsuitable for real-time AI interviews \cite{versatik2025webrtc}.

% ─────────────────────────────────────────────
\subsubsection{Primary Database \& Vector Search}
% ─────────────────────────────────────────────

The system requires ACID-compliant relational storage for user profiles, enrollment records, and session history, combined with high-throughput vector similarity search for course knowledge graph retrieval during peak concurrent usage.


\textbf{Decision: PostgreSQL with pgvector.} PostgreSQL with \texttt{pgvector} delivers 11.4$\times$ higher vector query throughput than Qdrant at 99\% recall \cite{tigerdata2025pgvector}, while simultaneously providing full ACID compliance for transactional data. The ability to combine vector similarity search with SQL JOINs in a single query eliminates the need for a separate vector database service, significantly reducing operational complexity. Pinecone's managed-only model introduces vendor lock-in and prevents on-premise deployment.

% ─────────────────────────────────────────────
\subsubsection{Graph Database}
% ─────────────────────────────────────────────

The Course Knowledge Graph requires efficient multi-hop traversal to identify concept relationships, prerequisite chains, and knowledge gaps for adaptive question generation. These traversal patterns are prohibitively expensive in purely relational models.


\textbf{Decision: Neo4j.} Neo4j's index-free adjacency model makes multi-hop traversals (e.g., ``find all prerequisite concepts of topic X within 3 hops'') orders of magnitude faster than ArangoDB's document-oriented storage layer \cite{neo4j2024graph}. The APOC plugin library provides ready-made graph algorithms (PageRank, community detection) that directly support knowledge gap analysis. ArangoDB's multi-model approach adds flexibility but sacrifices graph traversal performance compared to a purpose-built graph engine.

% ─────────────────────────────────────────────
\subsubsection{Caching \& Task Queue}
% ─────────────────────────────────────────────

The system requires a high-speed in-memory store for JWT session tokens, API rate limiting, and background task scheduling (SR card scan, quiz generation jobs). Offloading these workloads from PostgreSQL prevents write contention during peak examination periods.


\textbf{Decision: Redis.} Redis uniquely serves three roles in a single service: session cache, rate limiter, and task queue (via \texttt{arq}). Memcached matches Redis on raw throughput but lacks the data structures needed for rate limiting (sorted sets) and task queuing (streams/lists). RabbitMQ is purpose-built for message queuing but requires a separate caching solution, doubling operational overhead. Redis's optional persistence (AOF) also ensures queued jobs survive container restarts \cite{redis2024benchmark}.

% ─────────────────────────────────────────────
\subsubsection{Final Stack Summary}
% ─────────────────────────────────────────────

\begin{table}[H]
    \caption{Selected Technology Stack — aBridgeAI}
    \centering
    \renewcommand{\arraystretch}{1.3}
    \begin{tabular}{|p{0.20\textwidth}|p{0.20\textwidth}|p{0.50\textwidth}|}
        \hline
        \textbf{Component} & \textbf{Technology} & \textbf{Primary Justification} \\
        \hline
        Backend & FastAPI (Python) & Native AI/ML, async I/O, auto OpenAPI \\
        \hline
        Frontend & Vite + React + TanStack Query & Authenticated SPA — no SSR needed; 54\% smaller bundle vs Next.js, sub-50ms HMR \\
        \hline
        Real-Time & LiveKit & WebRTC UDP, built-in AI Agents SDK \\
        \hline
        Primary DB & PostgreSQL + pgvector & Full ACID + 11.4$\times$ vector throughput vs Qdrant \\
        \hline
        Graph DB & Neo4j & Native multi-hop traversal for KG \\
        \hline
        Cache / Queue & Redis & Session cache + rate limit + task queue in one \\
        \hline
        Object Storage & Garage (S3-compatible) & Self-hostable, no vendor lock-in \\
        \hline
    \end{tabular}
\end{table}
